\chapter{Limites et continuité}\label{ch:b1:continuity}

Le \omterm{thm:b1:continuity:ivt}{théorème des valeurs intermédiaires} et le \omterm{thm:b1:continuity:evt}{théorème des bornes
atteintes} ont été utilisés au lycée sur la foi du dessin. Avec les
suites (\cref{ch:b1:seq}) et la topologie de $\R$
(\cref{ch:b1:topology}) en main, ce chapitre les démontre --- et
complète la théorie par le théorème de la \omterm{def:b1:logic:inj}{bijection} monotone (qui
légitime $\arcsin$, $\operatorname{arcosh}$ et leurs semblables du
\cref{ch:b1:functions}) et le \omterm{thm:b1:continuity:heine}{théorème de Heine} sur la \omterm{def:b1:continuity:uniform}{continuité
uniforme}.

Dans tout le chapitre, $I$ est un \omterm{prop:b1:reals:intervals}{intervalle} et $f \colon I \to \R$~;
«~$x_0 \in \overline I$~» autorise les limites aux extrémités.

\section{Limites de fonctions}

\begin{definition}[Limite en un point]\label{def:b1:continuity:limit}
Soient $x_0 \in \overline{I}$ et $\ell \in \R$. On dit que $f(x) \to
\ell$ quand $x \to x_0$ lorsque\index{limite!d'une fonction}
\[
\forall \varepsilon > 0,\ \exists \delta > 0,\ \forall x \in I,
\qquad \abs{x - x_0} \leq \delta \implies \abs{f(x) - \ell} \leq
\varepsilon .
\]
Les limites en $\pm\infty$ et les limites infinies se définissent sur
le même modèle ($\abs{x - x_0} \leq \delta$ devient $x \geq M$~;
$\abs{f(x) - \ell} \leq \varepsilon$ devient $f(x) \geq M'$). Les
limites à droite et à gauche restreignent $x$ à $x > x_0$ (on écrit
alors $x \to x_0^+$) ou à $x < x_0$. La limite, lorsqu'elle existe,
est unique (même démonstration que la
\cref{prop:b1:seq:first}).
\end{definition}

\begin{example}[Une limite par encadrement]\label{ex:b1:continuity:floorlimit}
Calculons $\lim_{x \to 0} x\,\bigl\lfloor \frac1x \bigr\rfloor$.
L'encadrement de la \omterm{thm:b1:reals:floor}{partie entière} $\frac1x - 1 < \lfloor \frac1x
\rfloor \leq \frac1x$ donne, après multiplication par $x$ (attention
au signe~!)~:
\[
1 - x < x\Bigl\lfloor \frac1x \Bigr\rfloor \leq 1 \quad (x > 0),
\qquad
1 \leq x\Bigl\lfloor \frac1x \Bigr\rfloor < 1 - x \quad (x < 0),
\]
et les deux encadrements latéraux se referment sur $1$~: la limite
vaut $1$. Observons ce qui s'est passé~: $\lfloor \frac1x\rfloor$
seul fait des sauts sauvages près de $0$, mais le facteur $x$ dompte
chaque saut ($x$ fois un saut d'amplitude $1$ est petit), et seul
l'encadrement survit. L'idée à retenir~: les limites de produits d'un
facteur petit par un facteur à oscillation bornée relèvent de
l'encadrement, jamais du théorème d'opérations --- le théorème
d'opérations exige que les \emph{deux} facteurs convergent.
\end{example}

\begin{theorem}[Caractérisation séquentielle]\label{thm:b1:continuity:seqchar}
$f(x) \to \ell$ quand $x \to x_0$ si et seulement si~: pour
\emph{toute} suite $(u_n)$ de points de $I$ telle que $u_n \to x_0$,
on a $f(u_n) \to \ell$.
\end{theorem}

\begin{proof}
($\Rightarrow$) Soient $u_n \to x_0$ et $\varepsilon > 0$. Prenons
$\delta$ fourni par la définition, puis $N$ tel que $\abs{u_n - x_0}
\leq \delta$ pour $n \geq N$~: au-delà de $N$, $\abs{f(u_n) - \ell}
\leq \varepsilon$.

($\Leftarrow$) Par contraposée. Si $f \not\to \ell$~: un certain
$\varepsilon_0 > 0$ met en défaut tout $\delta$~; en choisissant
$\delta = \frac{1}{n+1}$, on produit $u_n \in I$ tel que $\abs{u_n -
x_0} \leq \frac{1}{n+1}$ et $\abs{f(u_n) - \ell} > \varepsilon_0$.
Alors $u_n \to x_0$ mais $f(u_n) \not\to \ell$.
\end{proof}

\begin{corollary}[Opérations, composition, ordre]\label{cor:b1:continuity:operations}
Sommes, produits, quotients (limite non nulle au dénominateur) de
limites se comportent comme pour les suites~; si $f \to \ell$ en $x_0$
et $g \to m$ en $\ell$, et si $g$ est définie autour de $\ell$ avec
$g(\ell) = m$, ou bien si $f \neq \ell$ au \omterm{def:b1:topology:open}{voisinage} de $x_0$, alors
$g \circ f \to m$ en $x_0$~; les limites conservent les inégalités
larges, et le \omterm{thm:b1:seq:order}{théorème d'encadrement} est valable.
\end{corollary}

\begin{proof}
Chaque \omterm{def:b1:logic:statement}{assertion} se transporte, via le
\cref{thm:b1:continuity:seqchar}, à son analogue pour les suites
(\cref{thm:b1:seq:operations,thm:b1:seq:order}). Pour la composition,
l'enchaînement direct mérite d'être écrit une fois~: soit
$\varepsilon > 0$~; la limite de $g$ en $\ell$ fournit $\eta > 0$ tel
que
\[
\abs{y - \ell} \leq \eta \implies \abs{g(y) - m} \leq
\varepsilon \quad (y \text{ dans le domaine de } g),
\]
où le cas $y = \ell$ est couvert parce que $g(\ell) = m$~; puis la
limite de $f$ en $x_0$ fournit $\delta > 0$ tel que $\abs{x -
x_0} \leq \delta \implies \abs{f(x) - \ell} \leq \eta$~; en
enchaînant, $\abs{x - x_0} \leq \delta$ donne $\abs{g(f(x)) - m} \leq
\varepsilon$. Dans l'hypothèse alternative ($f \neq \ell$ au
\omterm{def:b1:topology:open}{voisinage} de $x_0$), la valeur $y = \ell$ n'est jamais présentée à
$g$ et la valeur de $g$ en ce point est sans importance.
\end{proof}

\begin{example}[Pourquoi la clause de composition existe]\label{ex:b1:continuity:compositiontrap}
Soit $g(y) = 0$ pour $y \neq 0$ et $g(0) = 1$, et soit $f$
identiquement nulle. Alors $f(x) \to 0$ quand $x \to 0$, et $g(y)
\to 0$ quand $y \to 0$~; pourtant $g(f(x)) = g(0) = 1$ pour tout
$x$, donc $g \circ f \to 1 \neq 0$. La fonction intérieure est
posée \emph{exactement sur} la valeur interdite $\ell = 0$ pour
toujours, et la limite de $g$ en $0$ ignore ce que $g$ fait en
$0$. La clause du
\cref{cor:b1:continuity:operations} --- ou bien $g(\ell) = m$
(c'est-à-dire $g$ \omterm{def:b1:continuity:continuous}{continue} en $\ell$), ou bien $f \neq \ell$ au
\omterm{def:b1:topology:open}{voisinage} de $x_0$ --- est précisément ce qui exclut cela.
L'idée à retenir~: en pratique on compose des fonctions
\emph{\omterm{def:b1:continuity:continuous}{continues}} et la clause est gratuite~; elle ne mord que
lorsque les limites sont prises le long de \omterm{def:b1:topology:open}{voisinages} épointés,
raison pour laquelle la définition de $\lim_{x \to x_0}$ utilisée
dans ce livre inclut le point lorsqu'il appartient au domaine.
\end{example}

\section{Continuité}

\begin{definition}\label{def:b1:continuity:continuous}
$f$ est \emph{continue en $x_0 \in I$}\index{continuité} lorsque
$f(x) \to f(x_0)$ quand $x \to x_0$~; \emph{continue sur $I$}
lorsqu'elle est continue en tout point. D'après le
\cref{thm:b1:continuity:seqchar}~: $f$ est continue en $x_0$ si et
seulement si $f(u_n) \to f(x_0)$ pour toute suite $u_n \to x_0$ de
points de $I$.
\end{definition}

\begin{example}[Une taxonomie des discontinuités]\label{ex:b1:continuity:taxonomy}
Trois façons d'échouer en un point, par gravité croissante.
\emph{Éliminable}~: $f(x) = \frac{\sin x}{x}$ sur $\R^*$ a pour
limite $1$ en $0$~; poser $f(0) = 1$ répare tout --- la
discontinuité était un trou, pas un trait de caractère.
\emph{Saut}~: $\lfloor x \rfloor$ en un entier a des limites
latérales distinctes ($n - 1$ et $n$)~; aucun choix de valeur ne
peut les réconcilier, mais les deux demi-limites existent.
\emph{Essentielle}~: $\sin\frac1x$ en $0$ n'a aucune limite
latérale (\cref{exo:b1:continuity:1}) --- oscillation sans
apaisement. Les fonctions monotones ne peuvent produire que le
type intermédiaire (leurs limites latérales existent toujours,
comme \omterm{def:b1:reals:bounds}{bornes supérieures} et inférieures), et c'est pourquoi leur
\omterm{def:b1:logic:sets}{ensemble} de discontinuités est au plus dénombrable --- un
rationnel par saut. Les dérivées, d'après le théorème de Darboux
(\cref{exo:b1:derivative:10}), ne peuvent produire que le
dernier type~: une fonction présentant une discontinuité de
saut n'est jamais la dérivée de quoi que ce soit.
\end{example}

\begin{proposition}\label{prop:b1:continuity:algebra}
Sommes, produits, quotients (là où ils sont définis) et composées de
fonctions \omterm{def:b1:continuity:continuous}{continues} sont \omterm{def:b1:continuity:continuous}{continus}. Les \omterm{def:b1:poly:def}{polynômes}, les \omterm{def:b1:fractions:field}{fractions
rationnelles} (hors de leurs \omterm{def:b1:fractions:field}{pôles}), $\abs{\,\cdot\,}$, $\exp$, $\ln$,
les fonctions trigonométriques et hyperboliques et leurs réciproques
(\cref{ch:b1:functions}) sont \omterm{def:b1:continuity:continuous}{continus} sur leurs domaines.
\end{proposition}

\begin{proof}
Opérations~: le \cref{cor:b1:continuity:operations}. Les constantes
et l'identité sont \omterm{def:b1:continuity:continuous}{continues} directement d'après la définition
($\delta = \varepsilon$ convient pour l'identité, n'importe quel
$\delta$ pour les constantes)~; les produits de fonctions
\omterm{def:b1:continuity:continuous}{continues} étant \omterm{def:b1:continuity:continuous}{continus}, chaque monôme $a_k x^k$ suit par
récurrence sur $k$, et les sommes achèvent les \omterm{def:b1:poly:def}{polynômes}~; une
\omterm{def:b1:fractions:field}{fraction rationnelle} est un quotient de deux \omterm{def:b1:poly:def}{polynômes}, \omterm{def:b1:continuity:continuous}{continue}
partout où le dénominateur ne s'annule pas.
$\abs{\,\cdot\,}$~: seconde \omterm{prop:b1:complex:rules}{inégalité triangulaire},
$\bigl|\abs{f(x)} - \abs{f(x_0)}\bigr| \leq \abs{f(x) -
f(x_0)}$, donc le même $\delta$ fonctionne. Pour les fonctions
classiques, nous accordons ici la \omterm{def:b1:continuity:continuous}{continuité}~; la dérivabilité
(démontrée au \cref{ch:b1:derivative}) est plus forte.
\end{proof}

\begin{example}[Maximum et minimum de fonctions continues]\label{ex:b1:continuity:maxmin}
Si $f$ et $g$ sont \omterm{def:b1:continuity:continuous}{continues}, $\max(f, g)$ et $\min(f, g)$ le
sont aussi~: aucune disjonction de cas n'est nécessaire, grâce
aux identités
\[
\max(f, g) = \frac{f + g + \abs{f - g}}{2},
\qquad
\min(f, g) = \frac{f + g - \abs{f - g}}{2},
\]
et à la \omterm{def:b1:continuity:continuous}{continuité} des sommes et de $\abs{\,\cdot\,}$
(\cref{prop:b1:continuity:algebra}). En particulier $f^+ =
\max(f, 0)$ et $f^- = \max(-f, 0)$ sont \omterm{def:b1:continuity:continuous}{continues} avec $f =
f^+ - f^-$~: la décomposition selon le signe utilisée pour les
séries (\cref{ch:b1:series}) et, à pleine échelle, dans la
théorie de l'intégration du volume de Licence 3, ne coûte rien
en régularité.
\end{example}

\begin{theorem}[Théorème des valeurs intermédiaires]\label{thm:b1:continuity:ivt}
Soit $f$ \omterm{def:b1:continuity:continuous}{continue} sur $\intcc{a}{b}$ avec $f(a) \leq 0 \leq
f(b)$. Alors $f(c) = 0$ pour un certain $c \in \intcc{a}{b}$.
Par conséquent, une fonction \omterm{def:b1:continuity:continuous}{continue} sur un \omterm{prop:b1:reals:intervals}{intervalle} prend toute
valeur comprise entre deux de ses valeurs~: $f(I)$ est un
\omterm{prop:b1:reals:intervals}{intervalle}.\index{théorème des valeurs intermédiaires}
\end{theorem}

\begin{proof}
Dichotomie. Posons $a_0 = a$, $b_0 = b$. Étant donné $\intcc{a_n}{b_n}$
avec $f(a_n) \leq 0 \leq f(b_n)$, soit $m$ le milieu~: si $f(m) \leq
0$ on garde $\intcc{m}{b_n}$, sinon on garde $\intcc{a_n}{m}$~; les
conditions de signe persistent. Les suites $(a_n)$ et $(b_n)$ sont
adjacentes ($b_n - a_n = \frac{b-a}{2^n}$), de limite commune $c$
(\cref{thm:b1:seq:adjacent}). Par \omterm{def:b1:continuity:continuous}{continuité} et par le
\cref{thm:b1:seq:order}~: $f(c) = \lim f(a_n) \leq 0$ et $f(c) =
\lim f(b_n) \geq 0$, donc $f(c) = 0$.

Pour la conséquence~: étant données des valeurs $f(u) < v < f(w)$, on
applique ce qui précède à $x \mapsto f(x) - v$ sur le segment
d'extrémités $u$ et $w$~; ainsi $f(I)$ est convexe, c'est-à-dire un
\omterm{prop:b1:reals:intervals}{intervalle} (\cref{prop:b1:reals:intervals}).
\end{proof}

\begin{omfigure}
\begin{tikzpicture}
\begin{axis}[omaxis, width=8.4cm, height=5.6cm,
    xmin=0, xmax=4.4, ymin=-1.1, ymax=1.4,
    xtick={0.6,3.8}, xticklabels={$a$,$b$}, ytick=\empty]
  \addplot[omDef, very thick, domain=0.6:3.8, samples=160]
    {0.5*(x-2) + 0.4*sin(deg(2.5*x))};
  \draw[dashed, gray] (axis cs:0.6,0) -- (axis cs:0.6,-0.301);
  \draw[dashed, gray] (axis cs:3.8,0) -- (axis cs:3.8,0.870);
  \node[below, font=\small, omDef] at (axis cs:0.6,-0.32) {$f(a)$};
  \node[above, font=\small, omDef] at (axis cs:3.8,0.89) {$f(b)$};
\end{axis}
\end{tikzpicture}

{\small Une fonction \omterm{def:b1:continuity:continuous}{continue} vérifiant $f(a) < 0 < f(b)$ doit couper
l'axe~: la démonstration par dichotomie du
\cref{thm:b1:continuity:ivt} piège un point d'intersection entre deux
\omterm{thm:b1:seq:adjacent}{suites adjacentes}.}
\end{omfigure}

\begin{example}[Une équation, le protocole complet]\label{ex:b1:continuity:protocol}
Résolvons $\eu^x = 3 - x$ sur $\R$~: existence, unicité,
localisation. Posons $g(x) = \eu^x + x - 3$, \omterm{def:b1:continuity:continuous}{continue}.
Localisation et existence~: $g(0) = -2 < 0$ et $g(1) = \eu - 2 >
0$, donc le \omterm{thm:b1:continuity:ivt}{théorème des valeurs intermédiaires} plante une
solution dans $\intoo{0}{1}$. Unicité~: $g$ est la somme de
$\eu^x$, strictement croissante, et de $x - 3$, donc strictement
croissante sur $\R$~; une fonction strictement monotone prend
chaque valeur au plus une fois, donc la solution est unique sur
$\R$ tout entier (et pas seulement sur l'\omterm{prop:b1:reals:intervals}{intervalle} sondé). Le
protocole --- se ramener à $g = 0$, changement de signe pour
l'existence, monotonie pour l'unicité --- règle en trois lignes
la plupart des questions «~combien de solutions~?~», et les
tableaux de variations du \cref{ch:b1:derivative} l'étendent aux
$g$ non monotones en découpant $\R$ en branches monotones.
\end{example}

\begin{example}[La dichotomie comme algorithme]\label{ex:b1:continuity:dichotomy}
La démonstration du \cref{thm:b1:continuity:ivt} calcule. Prenons
$f(x) = x^3 + x - 1$~: $f(0) = -1 < 0 < 1 = f(1)$, donc une
racine est dans $\intoo{0}{1}$. En coupant en deux~:
\[
f(0.5) = -0.375 < 0, \qquad
f(0.75) = 0.171875 > 0, \qquad
f(0.625) = -0.130859375 < 0 ,
\]
de sorte que la racine est successivement piégée dans
$\intoo{0.5}{1}$, puis $\intoo{0.5}{0.75}$, puis
$\intoo{0.625}{0.75}$ (valeur exacte~: $c \approx 0.6823$). Après
$n$ étapes l'erreur est au plus $\frac{b - a}{2^n}$~: dix étapes
donnent trois décimales, vingt en donnent six. L'idée à retenir~:
le \omterm{thm:b1:continuity:ivt}{théorème des valeurs intermédiaires} n'est pas seulement un
énoncé d'existence --- sa démonstration par dichotomie est un
algorithme de recherche de racine garanti, quoique lent, auquel
il faut comparer la méthode de Newton du
\cref{ch:b1:derivative}, rapide mais locale.
\end{example}

\begin{theorem}[Théorème des bornes atteintes]\label{thm:b1:continuity:evt}
Une fonction \omterm{def:b1:continuity:continuous}{continue} sur un segment $\intcc{a}{b}$ est bornée et
atteint ses bornes~: il existe $c, d \in \intcc{a}{b}$ tels que
\[
f(c) = \inf_{\intcc{a}{b}} f,
\qquad
f(d) = \sup_{\intcc{a}{b}} f .
\]
Combiné au \cref{thm:b1:continuity:ivt}~: \emph{l'image \omterm{def:b1:continuity:continuous}{continue} d'un
segment est un segment} $\intcc{f(c)}{f(d)}$.
\index{théorème des bornes atteintes}
\end{theorem}

\begin{proof}
\emph{Majorée~:} sinon on choisit $u_n$ tel que $f(u_n) \geq n$. Par
compacité du segment (\cref{thm:b1:topology:compact}), une \omterm{def:b1:seq:subsequence}{suite
extraite} $u_{\varphi(n)} \to x \in \intcc{a}{b}$~; la \omterm{def:b1:continuity:continuous}{continuité} donne
$f(u_{\varphi(n)}) \to f(x)$, alors que $f(u_{\varphi(n)}) \geq
\varphi(n) \to +\infty$~: contradiction.

\emph{\omterm{def:b1:reals:bounds}{Borne supérieure} atteinte~:} soit $M = \sup f$ et choisissons
$v_n$ tel que $f(v_n) > M - \frac{1}{n+1}$
(\cref{prop:b1:reals:epsilon}). Extrayons $v_{\varphi(n)} \to d \in
\intcc{a}{b}$~: alors $f(d) = \lim f(v_{\varphi(n)}) = M$ par
encadrement. La \omterm{def:b1:reals:bounds}{borne inférieure} se traite avec $-f$.
\end{proof}

\begin{example}[Minimum strictement positif sur un segment]\label{ex:b1:continuity:posmin}
Soit $f$ \omterm{def:b1:continuity:continuous}{continue} sur $\intcc{0}{1}$ avec $f(x) > 0$ pour tout
$x$. Alors $\inf f > 0$~: par le \omterm{thm:b1:continuity:evt}{théorème des bornes atteintes},
la \omterm{def:b1:reals:bounds}{borne inférieure} est une \emph{valeur} $f(c)$, et $f(c) > 0$
par hypothèse. Ainsi une fonction \omterm{def:b1:continuity:continuous}{continue} strictement positive
sur un segment est minorée par une constante strictement
positive --- un argument de deux lignes utilisé une douzaine de
fois dans les chapitres à venir (dénominateurs sous contrôle,
encadrements par des fonctions en escalier, majorations
d'erreur). Sur un \omterm{prop:b1:reals:intervals}{intervalle} non compact, cela échoue
spectaculairement~: $f(x) = x$ sur $\intoc{0}{1}$ est \omterm{def:b1:continuity:continuous}{continue} et
strictement positive avec $\inf f = 0$, non atteint. L'idée à
retenir~: «~strictement positive~» se hausse en «~uniformément
positive~» exactement lorsque le domaine est compact~; chaque
hypothèse du \omterm{thm:b1:continuity:evt}{théorème des bornes atteintes} porte le toit.
\end{example}

\begin{remark}[Pièges classiques autour des trois théorèmes]
(i) \emph{Images \omterm{def:b1:continuity:continuous}{continues}}~: seuls les \emph{segments} sont
robustes. L'image \omterm{def:b1:continuity:continuous}{continue} d'un \omterm{prop:b1:reals:intervals}{intervalle} \omterm{def:b1:topology:open}{ouvert} n'est pas
nécessairement \omterm{def:b1:topology:open}{ouverte} ($x \mapsto x^2$ envoie $\intoo{-1}{1}$
sur $\intco{0}{1}$), l'image d'un \omterm{def:b1:topology:closed}{fermé} n'est pas nécessairement
\omterm{def:b1:topology:closed}{fermée} ($\arctan$ envoie le \omterm{def:b1:topology:closed}{fermé} $\R$ sur l'\omterm{def:b1:topology:open}{ouvert}
$\intoo{-\frac\pi2}{\frac\pi2}$)~; mais l'image d'un segment est
un segment (\cref{thm:b1:continuity:evt}). (ii) \emph{Le
\omterm{thm:b1:continuity:ivt}{théorème des valeurs intermédiaires} exige un \omterm{prop:b1:reals:intervals}{intervalle}}~: la
fonction $x \mapsto \frac1x$, \omterm{def:b1:continuity:continuous}{continue} sur $\intco{-1}{0} \cup
\intoc{0}{1}$, prend les valeurs $-1$ et $1$ sans jamais
s'annuler --- son domaine est fait de deux morceaux disjoints, et
la valeur $0$ tombe dans le trou~; nommez toujours l'\omterm{prop:b1:reals:intervals}{intervalle}
sur lequel le théorème est appliqué. (iii) \emph{La \omterm{def:b1:continuity:uniform}{continuité
uniforme} est une propriété du couple (fonction, \omterm{def:b1:logic:sets}{ensemble})}~:
$x^2$ est \omterm{def:b1:continuity:uniform}{uniformément continue} sur tout segment mais pas sur
$\R$ (\cref{ex:b1:continuity:notuniform}) --- les mots
«~\omterm{def:b1:continuity:uniform}{uniformément continue}~» sans domaine n'ont aucun sens. (iv)
\emph{La \omterm{def:b1:continuity:continuous}{continuité} de la réciproque n'est pas formelle}~: elle
vaut sur les \omterm{prop:b1:reals:intervals}{intervalles} grâce à la monotonie
(\cref{thm:b1:continuity:bijection}), mais une \omterm{def:b1:logic:inj}{bijection} \omterm{def:b1:continuity:continuous}{continue}
entre réunions d'\omterm{prop:b1:reals:intervals}{intervalles} peut avoir une réciproque
discontinue --- la remarque qui suit ce théorème est là parce que
les étudiants le citent sans l'hypothèse d'\omterm{prop:b1:reals:intervals}{intervalle}.
\end{remark}

\section{Fonctions monotones et fonctions réciproques}

\begin{theorem}[Théorème de la bijection monotone]\label{thm:b1:continuity:bijection}
Soit $f$ \omterm{def:b1:continuity:continuous}{continue} et strictement monotone sur un \omterm{prop:b1:reals:intervals}{intervalle} $I$. Alors
\begin{enumerate}
  \item $f$ est une \omterm{def:b1:logic:inj}{bijection} de $I$ sur l'\omterm{prop:b1:reals:intervals}{intervalle} $J = f(I)$~;
  \item la réciproque $f^{-1} \colon J \to I$ est strictement monotone
        (de même sens) et \emph{\omterm{def:b1:continuity:continuous}{continue}}.
\end{enumerate}
\end{theorem}

\begin{proof}
Supposons $f$ strictement croissante. (1) L'\omterm{def:b1:logic:inj}{injectivité} est immédiate
par stricte monotonie~; la \omterm{def:b1:logic:inj}{surjectivité} sur $f(I)$ est triviale, et
$f(I)$ est un \omterm{prop:b1:reals:intervals}{intervalle} d'après le \cref{thm:b1:continuity:ivt}.

(2) $f^{-1}$ est strictement croissante~: si $y < y'$ dans $J$ mais
$f^{-1}(y) \geq f^{-1}(y')$, appliquer $f$, croissante, donne $y \geq
y'$, absurde. \omterm{def:b1:continuity:continuous}{Continuité} de $f^{-1}$ en $y_0 = f(x_0) \in J$~: soit
$\varepsilon > 0$. Supposons d'abord $x_0$ \omterm{def:b1:topology:closure}{intérieur} à $I$, et
diminuons $\varepsilon$ de sorte que $x_0 \pm \varepsilon \in I$~:
leurs images vérifient $f(x_0 - \varepsilon) < y_0 < f(x_0 +
\varepsilon)$. Prenons
$\delta = \min\bigl(y_0 - f(x_0 - \varepsilon),\, f(x_0 +
\varepsilon) - y_0\bigr) > 0$~: pour $\abs{y - y_0} \leq \delta$, la
monotonie de $f^{-1}$ coince $f^{-1}(y)$ entre $x_0 -
\varepsilon$ et $x_0 + \varepsilon$. Si $x_0$ est, disons,
l'extrémité gauche de $I$, seul $x_0 + \varepsilon$ est disponible~:
alors $y_0 = \min J$ (monotonie), tout $y \in J$ tel que $y - y_0
\leq f(x_0 + \varepsilon) - y_0$ vérifie $x_0 \leq f^{-1}(y)
\leq x_0 + \varepsilon$, et l'estimation unilatérale est exactement
la \omterm{def:b1:continuity:continuous}{continuité} en une extrémité~; l'extrémité droite est symétrique.
(Remarquons que la \omterm{def:b1:continuity:continuous}{continuité} de $f^{-1}$
n'a \emph{pas} été déduite de celle de $f$ par symétrie --- c'est la
monotonie sur un \omterm{prop:b1:reals:intervals}{intervalle} qui fait le travail.)
\end{proof}

\begin{remark}
Ce théorème est ce que la \cref{def:b1:functions:arc} et la
\cref{prop:b1:functions:invhyp} ont utilisé en silence~: $\arcsin$,
$\arccos$, $\arctan$, $\operatorname{arsinh}$, \dots\ sont \omterm{def:b1:continuity:continuous}{continues}.
Un complément (\cref{exo:b1:continuity:10})~: une fonction \omterm{def:b1:continuity:continuous}{continue}
\emph{\omterm{def:b1:logic:inj}{injective}} sur un \omterm{prop:b1:reals:intervals}{intervalle} est automatiquement strictement
monotone, de sorte que l'hypothèse de monotonie ne coûte rien.
\end{remark}

\begin{example}[Racines de tous les ordres]\label{ex:b1:continuity:nthroot}
Pour $n \in \N^*$, la fonction $f(x) = x^n$ est \omterm{def:b1:continuity:continuous}{continue} et
strictement croissante sur $\intco{0}{+\infty}$, avec $f(0) = 0$
et $f(x) \to +\infty$~: son image est $\intco{0}{+\infty}$
(\cref{thm:b1:continuity:ivt} pour la structure d'\omterm{prop:b1:reals:intervals}{intervalle}). Le
théorème de la \omterm{def:b1:logic:inj}{bijection} monotone livre alors, d'un seul coup,
une réciproque strictement croissante et \emph{\omterm{def:b1:continuity:continuous}{continue}}
\[
x \mapsto x^{1/n} \colon \intco{0}{+\infty} \to
\intco{0}{+\infty} :
\]
existence, unicité et \omterm{def:b1:continuity:continuous}{continuité} des racines $n$-ièmes, sans
le moindre calcul. À comparer avec l'\cref{exo:b1:reals:12}, qui
construisait $\sqrt y$ à la main à partir de la \omterm{def:b1:reals:bounds}{borne
supérieure}~: un chapitre de théorie a comprimé cette page de
travail en deux lignes, et ces deux mêmes lignes avaient déjà
légitimé $\arcsin$, $\arctan$ et $\operatorname{arsinh}$. L'idée
à retenir~: un bon théorème est du travail mis en réserve.
\end{example}

\section{Continuité uniforme}

\begin{definition}\label{def:b1:continuity:uniform}
$f \colon I \to \R$ est \emph{uniformément
continue}\index{continuité uniforme} lorsque
\[
\forall \varepsilon > 0,\ \exists \delta > 0,\ \forall x, y \in I,
\qquad \abs{x - y} \leq \delta \implies \abs{f(x) - f(y)} \leq
\varepsilon .
\]
Le point important~: $\delta$ ne dépend que de $\varepsilon$, et pas
de l'endroit où l'on se trouve dans $I$. La \omterm{def:b1:continuity:continuous}{continuité} uniforme
entraîne la \omterm{def:b1:continuity:continuous}{continuité}~; une fonction \emph{lipschitzienne}
($\abs{f(x) - f(y)} \leq k \abs{x - y}$) est uniformément \omterm{def:b1:continuity:continuous}{continue}
($\delta = \varepsilon/k$).
\end{definition}

\begin{example}\label{ex:b1:continuity:notuniform}
$x \mapsto x^2$ est \omterm{def:b1:continuity:continuous}{continue} sur $\R$ mais \emph{pas} \omterm{def:b1:continuity:uniform}{uniformément
continue}~: $\abs{(n + \frac 1n)^2 - n^2} = 2 + \frac{1}{n^2} \geq 2$
alors que les arguments sont à distance $\frac 1n \to 0$. Sur tout
\omterm{prop:b1:reals:intervals}{intervalle} borné elle est lipschitzienne, donc \omterm{def:b1:continuity:uniform}{uniformément continue}
--- en accord avec le \omterm{thm:b1:continuity:heine}{théorème de Heine} ci-dessous.
\end{example}

\begin{example}[Modules d'uniformité, explicitement]\label{ex:b1:continuity:moduli}
Sur un segment, Heine garantit un $\delta$ uniforme~; souvent on
peut aussi le \emph{calculer}. Pour $f(x) = x^2$ sur
$\intcc{0}{10}$~:
\[
\abs{x^2 - y^2} = \abs{x + y}\,\abs{x - y} \leq 20\,\abs{x - y},
\]
donc $\delta = \frac{\varepsilon}{20}$ convient uniformément (un
module lipschitzien, linéaire en $\varepsilon$). Pour $\sqrt x$
sur $\intcc{0}{1}$~: $\abs{\sqrt x - \sqrt y} \leq \sqrt{\abs{x -
y}}$ (\cref{exo:b1:continuity:9}), donc $\delta = \varepsilon^2$
convient --- uniforme mais \emph{pas} linéaire~: près de $0$ la
racine carrée est raide, et le prix apparaît dans l'exposant de
$\varepsilon$, non dans un défaut d'uniformité. L'idée à
retenir~: la \omterm{def:b1:continuity:uniform}{continuité uniforme} est un dégradé, non un oui/non
--- la fonction $\delta(\varepsilon)$, appelée \omterm{def:b1:complex:field}{module}, mesure le
coût de l'uniformité, et le caractère lipschitzien en est
simplement la meilleure note.
\end{example}

\begin{theorem}[Heine]\label{thm:b1:continuity:heine}
Une fonction \omterm{def:b1:continuity:continuous}{continue} sur un \emph{segment} $\intcc{a}{b}$ est
\omterm{def:b1:continuity:uniform}{uniformément continue}.\index{théorème de Heine}
\end{theorem}

\begin{proof}
Par l'absurde~: supposons qu'un certain $\varepsilon_0 > 0$ mette en
défaut tout $\delta$. Avec $\delta = \frac{1}{n+1}$, choisissons
$x_n, y_n \in \intcc{a}{b}$ tels que $\abs{x_n - y_n} \leq
\frac{1}{n+1}$ et $\abs{f(x_n) - f(y_n)} > \varepsilon_0$. La
compacité (\cref{thm:b1:topology:compact}) extrait $x_{\varphi(n)} \to
c \in \intcc{a}{b}$~; alors $y_{\varphi(n)} \to c$ également
(encadrement sur $\abs{x - y}$). La \omterm{def:b1:continuity:continuous}{continuité} en $c$ donne
$f(x_{\varphi(n)}) \to f(c)$ et $f(y_{\varphi(n)}) \to f(c)$, donc
$\abs{f(x_{\varphi(n)}) - f(y_{\varphi(n)})} \to 0 < \varepsilon_0$~:
contradiction.
\end{proof}

\begin{example}[Bornée, continue, et pourtant pas uniformément]\label{ex:b1:continuity:sinx2}
La fonction $f(x) = \sin(x^2)$ est \omterm{def:b1:continuity:continuous}{continue} et bornée sur $\R$,
mais pas \omterm{def:b1:continuity:uniform}{uniformément continue}. Prenons
\[
x_n = \sqrt{2\pi n}, \qquad y_n = \sqrt{2\pi n + \tfrac\pi2}:
\qquad
y_n - x_n = \frac{\pi/2}{\sqrt{2\pi n + \frac\pi2} +
\sqrt{2\pi n}} \longrightarrow 0 ,
\]
et pourtant $f(y_n) - f(x_n) = \sin\bigl(2\pi n + \frac\pi2\bigr)
- \sin(2\pi n) = 1 - 0 = 1$ pour tout $n$~: aucun $\delta$ unique
ne peut servir $\varepsilon = \frac12$ partout. Géométriquement,
les oscillations de $\sin(x^2)$ \emph{accélèrent}~: le graphe
accomplit une onde complète sur des fenêtres de plus en plus
courtes, de sorte que l'échelle horizontale exigée par un
$\varepsilon$ donné tend vers zéro quand $x$ grandit. L'idée à
retenir~: le caractère borné n'achète pas l'uniformité (cet
exemple), et le caractère non borné ne l'interdit pas ($\sqrt x$,
\cref{exo:b1:continuity:9})~; ce qui décide, c'est le
\emph{module d'oscillation}, et le \omterm{thm:b1:continuity:heine}{théorème de Heine} dit que les
domaines compacts le disciplinent automatiquement.
\end{example}

\begin{example}[L'expérience de la calculatrice, expliquée]\label{ex:b1:continuity:cositeration}
Tapez un nombre quelconque sur une calculatrice et appuyez
plusieurs fois sur $\cos$~: l'affichage se verrouille sur
$0.7390851\dots$ Pourquoi~? Après une pression la valeur est dans
$\intcc{-1}{1}$, après deux dans $\intcc{\cos 1}{1}
\subseteq \intcc{0.54}{1}$, \omterm{prop:b1:reals:intervals}{intervalle} stable pour $\cos$. Sur
celui-ci, $\abs{\cos a - \cos b} \leq \sin(1)\,\abs{a - b}$ avec
$\sin 1 = 0.841\dots < 1$ (la majoration par produit-somme du
\cref{pb:b1:seq:1}, ou l'inégalité des accroissements finis du
\cref{ch:b1:derivative})~: l'itération est contractante, donc par
l'étape de contrôle de l'erreur de la \cref{met:b1:seq:recurrent},
\[
\abs{u_n - c} \leq (0.842)^{\,n-2}\,\abs{u_2 - c}
\longrightarrow 0 ,
\]
où $c$ est l'unique point fixe $\cos c = c$
(\cref{exo:b1:continuity:6}). Une quarantaine de pressions
achètent trois décimales ($0.842^{40} \approx 10^{-3}$) --- une
vitesse géométrique, plus douce que le $2^{-n}$ de la dichotomie,
mais chaque pression coûte une touche là où chaque étape de
dichotomie coûte une évaluation de signe complète. L'idée à
retenir~: l'image du point fixe du \cref{ch:b1:seq} et les
théorèmes d'existence de ce chapitre sont les deux moitiés d'une
même histoire --- le \omterm{thm:b1:continuity:ivt}{théorème des valeurs intermédiaires} trouve
$c$, la contraction l'atteint.
\end{example}

\begin{remark}[Où ces théorèmes serviront ensuite]
Les trois piliers de ce chapitre alimentent chacun un chapitre
ultérieur. Le \omterm{thm:b1:continuity:ivt}{théorème des valeurs intermédiaires} nourrit tout
argument d'existence de solutions ainsi que le théorème de la
\omterm{def:b1:logic:inj}{bijection} monotone~; le \omterm{thm:b1:continuity:evt}{théorème des bornes atteintes} transforme
les problèmes d'optimisation en théorèmes (Rolle et le théorème
des accroissements finis du \cref{ch:b1:derivative} partent
exactement de là)~; le \omterm{thm:b1:continuity:heine}{théorème de Heine} est la raison pour
laquelle les fonctions \omterm{def:b1:continuity:continuous}{continues} sur un segment sont intégrables
au \cref{ch:b1:integration} --- le $\delta$ uniforme est ce qui
fait converger les sommes de Riemann. Dans le volume de Licence
2, le même trio reparaît dans les espaces vectoriels normés, la
compacité y faisant le travail que les segments font ici.
\end{remark}

\begin{remark}[Perspectives dans ce volume]
La \omterm{def:b1:continuity:continuous}{continuité} va être dépassée sans jamais être mise à la
retraite. Le \cref{ch:b1:derivative} la renforce en dérivabilité
et lui rend la politesse (dérivable entraîne \omterm{def:b1:continuity:continuous}{continue})~; le
\cref{ch:b1:integration} repose deux fois sur elle, par Heine
pour la construction et par le théorème fondamental, dont l'objet
central $x \mapsto \int_a^x f$ hausse une $f$ simplement \omterm{def:b1:continuity:continuous}{continue}
en une primitive de classe $C^1$. Au
\cref{ch:b1:multivar}, la \omterm{def:b1:continuity:continuous}{continuité} des fonctions de deux
variables recèle un piège qui mérite d'être annoncé~: la fonction $\frac{xy}{x^2 +
y^2}$ (prolongée par $0$) est \omterm{def:b1:continuity:continuous}{continue} en $x$ pour chaque $y$
fixé et en $y$ pour chaque $x$ fixé, et pourtant elle n'est pas
\omterm{def:b1:continuity:continuous}{continue} à l'origine --- le long de la diagonale $x = y$ elle vaut
constamment $\frac12$. La \omterm{def:b1:continuity:continuous}{continuité} séparée est strictement plus
faible que la \omterm{def:b1:continuity:continuous}{continuité}~: la caractérisation séquentielle
survit au passage à $\R^2$, mais il faut autoriser les suites à
s'approcher depuis \emph{toutes} les directions, et pas seulement
le long des axes.
\end{remark}

\section{Exercices}

\begin{exercise}[$\star$]\label{exo:b1:continuity:1}
À l'aide de la caractérisation séquentielle, démontrer que $x \mapsto
\sin\frac 1x$ n'a pas de limite en $0^+$ \emph{(exhiber deux
suites)}. La fonction $x \mapsto x \sin\frac 1x$ en a-t-elle une~?
\end{exercise}

\begin{exercise}[$\star$]\label{exo:b1:continuity:2}
Étudier la \omterm{def:b1:continuity:continuous}{continuité} sur $\R$ de $f(x) = \lfloor x \rfloor$, de
$g(x) = x - \lfloor x \rfloor$, et de $h(x) = \lfloor x \rfloor +
(x - \lfloor x\rfloor)^2$.
\end{exercise}

\begin{exercise}[$\star$]\label{exo:b1:continuity:3}
Démontrer que l'équation $x^5 - 3x + 1 = 0$ a au moins trois
solutions réelles \emph{(évaluer en des points bien choisis et
appliquer le \cref{thm:b1:continuity:ivt} sur trois segments
disjoints)}.
\end{exercise}

\begin{exercise}[$\star$]\label{exo:b1:continuity:4}
Démontrer que tout \omterm{def:b1:poly:def}{polynôme} de degré impair a une racine réelle.
\end{exercise}

\begin{exercise}[$\star\star$]\label{exo:b1:continuity:5}
(Point fixe) Soit $f \colon \intcc{0}{1} \to \intcc{0}{1}$ \omterm{def:b1:continuity:continuous}{continue}.
Démontrer que $f$ a un point fixe~: $f(c) = c$ pour un certain $c$.
Illustrer qu'on ne peut renoncer ni à la \omterm{def:b1:continuity:continuous}{continuité} ni au segment.
\end{exercise}

\begin{exercise}[$\star\star$]\label{exo:b1:continuity:6}
Démontrer que l'équation $\cos x = x$ a exactement une solution
réelle, et qu'elle appartient à $\intoo{0}{1}$.
\end{exercise}

\begin{exercise}[$\star\star$]\label{exo:b1:continuity:7}
Soit $f \colon \R \to \R$ \omterm{def:b1:continuity:continuous}{continue} avec $f(x) \to +\infty$ quand $x
\to \pm\infty$. Démontrer que $f$ atteint un minimum global sur $\R$.
\emph{(Se ramener à un segment contenant un \omterm{def:b1:logic:sets}{ensemble} de sous-niveau.)}
\end{exercise}

\begin{exercise}[$\star\star$]\label{exo:b1:continuity:8}
Soit $f \colon \R \to \R$ \omterm{def:b1:continuity:continuous}{continue} et périodique (de période $T >
0$). Démontrer que $f$ est bornée et atteint ses bornes, et qu'il
existe $c$ tel que $f(c + \frac T2) = f(c)$. \emph{(Pour le second
point, étudier $g(x) = f(x + \frac T2) - f(x)$ sur une période.)}
\end{exercise}

\begin{exercise}[$\star\star$]\label{exo:b1:continuity:9}
Démontrer que $x \mapsto \sqrt x$ est \omterm{def:b1:continuity:uniform}{uniformément continue} sur
$\intco{0}{+\infty}$, bien qu'elle ne soit pas lipschitzienne près de
$0$. \emph{(Démontrer et utiliser $\abs{\sqrt x - \sqrt y} \leq
\sqrt{\abs{x - y}}$.)}
\end{exercise}

\begin{exercise}[$\star\star\star$]\label{exo:b1:continuity:10}
Soit $f$ \omterm{def:b1:continuity:continuous}{continue} et \omterm{def:b1:logic:inj}{injective} sur un \omterm{prop:b1:reals:intervals}{intervalle} $I$. Démontrer que
$f$ est strictement monotone. \emph{Indication~: sinon, il existe $a
< b < c$ avec, disons, $f(b) > f(a)$ et $f(b) > f(c)$~; appliquer le
\omterm{thm:b1:continuity:ivt}{théorème des valeurs intermédiaires} à une valeur comprise entre
$\max(f(a), f(c))$ et $f(b)$ de part et d'autre de $b$.}
\end{exercise}

\begin{exercise}[$\star\star\star$]\label{exo:b1:continuity:11}
(\omterm{pb:b1:continuity:1}{Équation fonctionnelle de Cauchy}, cas \omterm{def:b1:continuity:continuous}{continu}) Soit $f \colon \R \to
\R$ \omterm{def:b1:continuity:continuous}{continue} avec $f(x + y) = f(x) + f(y)$ pour tous $x, y$.
Démontrer que $f(x) = f(1)\,x$ pour tout $x$~: d'abord sur $\N$,
$\Z$, $\Q$ (par la seule additivité), puis sur $\R$ par \omterm{def:b1:continuity:continuous}{continuité} et
\omterm{def:b1:topology:dense}{densité} (\cref{thm:b1:reals:density}).
\end{exercise}

\begin{exercise}[$\star\star\star$]\label{exo:b1:continuity:12}
Soit $f \colon \intco{0}{+\infty} \to \R$ \omterm{def:b1:continuity:continuous}{continue} avec
$f(x) \to \ell \in \R$ quand $x \to +\infty$. Démontrer que $f$ est
\emph{uniformément} \omterm{def:b1:continuity:continuous}{continue} sur $\intco{0}{+\infty}$. \emph{(Couper
en un $A$ grand~: Heine sur $\intcc{0}{A+1}$, la limite au-delà de
$A$~; faire se recouvrir les deux régimes.)}
\end{exercise}

\section{Problème~: l'équation fonctionnelle de Cauchy et ses sœurs}

\begin{problem}[{Devoir maison --- $f(x + y) = f(x) + f(y)$~:
la régularité force la linéarité, et le portrait des monstres}]
\label{pb:b1:continuity:1}
Quelles fonctions vérifient $f(x + y) = f(x) + f(y)$ pour tous
réels $x, y$~? Cauchy a posé la question en 1821~; la réponse est
un paradigme. L'\cref{exo:b1:continuity:11} montre qu'une telle
fonction \emph{additive} est linéaire sur $\Q$ et que la
\omterm{def:b1:continuity:continuous}{continuité} complète force $f(x) = cx$. Ce problème affaiblit
l'hypothèse de façon spectaculaire --- la \omterm{def:b1:continuity:continuous}{continuité} en un
\emph{seul} point, ou la monotonie, ou le simple caractère borné
sur un petit \omterm{prop:b1:reals:intervals}{intervalle}, chacun suffit --- puis peint le portrait
d'une hypothétique solution non linéaire (son graphe remplit le
plan), résout les équations sœurs qui caractérisent $\eu^{cx}$,
$c\ln x$, $x^c$ et $cx^2$, et se termine par l'équation de Jensen
et le théorème \emph{convexe au milieu $+$ \omterm{def:b1:continuity:continuous}{continue} $\implies$
convexe}.\index{équation fonctionnelle de
Cauchy}\index{équation fonctionnelle} Dans tout le problème,
\emph{additive} signifie~: $f(x + y) = f(x) + f(y)$ pour tous $x,
y \in \R$.

\medskip
\noindent\textbf{Partie I --- $\Q$-linéarité, et un seul point de
\omterm{def:b1:continuity:continuous}{continuité}.}
\begin{enumerate}
  \item Soit $f$ additive. D'après l'\cref{exo:b1:continuity:11},
        $f(r) = f(1)\,r$ pour $r$ rationnel. Démontrer l'énoncé
        plus fin utilisé ci-dessous~: pour \emph{tout} $x \in \R$
        et tout $r \in \Q$, $f(rx) = r\,f(x)$ ($f$ est
        $\Q$-\emph{linéaire}).
  \item Supposons la fonction additive $f$ \omterm{def:b1:continuity:continuous}{continue} en un seul
        point $x_0$. Montrer que $f$ est \omterm{def:b1:continuity:continuous}{continue} partout
        \emph{(calculer $f(x + h) - f(x)$ en fonction de $f(x_0 +
        h) - f(x_0)$)}, donc que $f(x) = f(1)\,x$.
  \item Montrer qu'une fonction additive est déterminée par sa
        restriction à tout \omterm{def:b1:structures:subgroup}{sous-groupe} \omterm{def:b1:topology:dense}{dense}~: si deux fonctions
        additives coïncident sur $\Z + \sqrt2\,\Z$ et sont toutes
        deux \omterm{def:b1:continuity:continuous}{continues}, elles sont égales --- alors que, sans
        \omterm{def:b1:continuity:continuous}{continuité}, prescrire $f(1) = 0$ et $f(\sqrt 2) = 1$ est
        compatible avec la $\Q$-linéarité sur le \omterm{def:b1:structures:subgroup}{sous-groupe}.
        Calculer $f(m + n\sqrt2)$ pour cette prescription.
  \item Soit $f$ additive et majorée par $M$ sur un certain
        \omterm{prop:b1:reals:intervals}{intervalle} $\intcc{a}{b}$ avec $a < b$. Montrer que $f$
        est majorée sur $\intcc{0}{\ell}$, $\ell = b - a$
        \emph{(translater de $a$)}.
\end{enumerate}

\medskip
\noindent\textbf{Partie II --- L'échelle de régularité.}
\begin{enumerate}[resume]
  \item Suite de la question 4~: en utilisant $f(\ell - t) + f(t) =
        f(\ell)$, montrer que $f$ est aussi \emph{minorée} sur
        $\intcc{0}{\ell}$~: $\abs f \leq C$ sur cet \omterm{prop:b1:reals:intervals}{intervalle}.
  \item Montrer que $\abs{f(t)} \leq \frac{C}{n}$ pour $t \in
        \intcc{0}{\ell/n}$, et en déduire que $f$ est \omterm{def:b1:continuity:continuous}{continue} en
        $0$ \emph{(l'imparité règle le côté gauche)}, donc partout
        (question 2)~: une fonction additive bornée sur un
        \omterm{prop:b1:reals:intervals}{intervalle} est linéaire.
  \item En déduire le cas monotone~: une fonction additive
        croissante sur un certain $\intcc{a}{b}$ ($a < b$) est
        $f(x) = cx$ avec $c \geq 0$.
  \item Assembler l'\emph{échelle de régularité}~: pour $f$
        additive, les \omterm{def:b1:logic:statement}{assertions} suivantes sont équivalentes ---
        (a) $f(x) = cx$~; (b) $f$ \omterm{def:b1:continuity:continuous}{continue}~; (c) $f$ \omterm{def:b1:continuity:continuous}{continue} en
        un point~; (d) $f$ monotone sur un \omterm{prop:b1:reals:intervals}{intervalle} non
        dégénéré~; (e) $f$ bornée sur un \omterm{prop:b1:reals:intervals}{intervalle} non dégénéré.
        Organiser les implications de sorte que chacune soit ou
        bien triviale, ou bien déjà démontrée.
  \item Vérifier que les questions 4 à 6 n'ont consommé qu'une
        majoration \emph{par en haut}~: une fonction additive
        majorée sur un \omterm{prop:b1:reals:intervals}{intervalle} non dégénéré est déjà linéaire.
        En déduire l'énoncé miroir pour une minoration, et
        consigner la forme la plus forte ainsi obtenue du barreau
        (e) de l'échelle.
\end{enumerate}

\medskip
\noindent\textbf{Partie III --- Portrait d'un monstre.} Supposons
maintenant $f$ additive mais \emph{non} linéaire.
\begin{enumerate}[resume]
  \item Montrer qu'il existe des réels non nuls $u, v$ tels que
        $\dfrac{f(u)}{u} \neq \dfrac{f(v)}{v}$, et que les
        vecteurs $(u, f(u))$ et $(v, f(v))$ engendrent le plan
        (leur déterminant $u f(v) - v f(u)$ est non nul).
  \item Montrer que le graphe de $f$ contient tous les points
        \[
        r\,(u, f(u)) + s\,(v, f(v)), \qquad r, s \in \Q ,
        \]
        et en déduire que le graphe est \emph{\omterm{def:b1:topology:dense}{dense} dans $\R^2$}~:
        pour tout point $(x_0, y_0)$ du plan et tout
        $\varepsilon > 0$, un certain $(x, f(x))$ en est à
        distance moindre que $\varepsilon$ \emph{(résoudre le
        système réel $2\times2$, puis approcher les coefficients
        réels par des rationnels)}.
  \item Déduire de la question 11 le portrait complet~: une
        fonction additive non linéaire est non bornée sur tout
        \omterm{prop:b1:reals:intervals}{intervalle} non dégénéré, discontinue en tout point,
        monotone sur aucun \omterm{prop:b1:reals:intervals}{intervalle}, et l'image qu'elle donne
        de tout \omterm{prop:b1:reals:intervals}{intervalle} est \omterm{def:b1:topology:dense}{dense} dans $\R$. Concilier avec la
        question 8.
  \item Les monstres existent --- sur un \omterm{def:b1:structures:subgroup}{sous-groupe} \omterm{def:b1:topology:dense}{dense}, de
        façon constructive~: sur $G = \Z + \sqrt2\,\Z$, posons
        $f(m + n\sqrt2) = n$. Montrer que $f$ est bien définie et
        additive sur $G$, et que $f$ est non bornée sur $G \cap
        \intoo{0}{\varepsilon}$ pour tout $\varepsilon > 0$
        \emph{(à $N$ fixé, seuls un nombre fini de $g = m +
        n\sqrt2 \in \intoo{0}{1}$ vérifient $\abs n \leq N$~; or
        $G \cap \intoo{0}{\varepsilon}$ est infini)}. Expliquer
        en un paragraphe pourquoi prolonger un tel $f$ à $\R$ tout
        entier exige une base de $\R$ comme $\Q$-espace vectoriel
        (une \emph{base de Hamel}), dont l'existence relève de
        l'axiome du choix, au-delà de ce volume.
\end{enumerate}

\medskip
\noindent\textbf{Partie IV --- Les équations sœurs.} Toutes les
fonctions sont ici \omterm{def:b1:continuity:continuous}{continues}.
\begin{enumerate}[resume]
  \item Soit $f \colon \R \to \R$ \omterm{def:b1:continuity:continuous}{continue}, non identiquement
        nulle, avec $f(x + y) = f(x)f(y)$. Montrer que $f(x) =
        f\bigl(\frac x2\bigr)^2 \geq 0$, puis que $f > 0$ partout,
        puis que $f(x) = \eu^{cx}$ pour un certain $c$~: les
        exponentielles sont exactement les \omterm{def:b1:structures:morphism}{morphismes} \omterm{def:b1:continuity:continuous}{continus} de
        $(\R, +)$ dans $(\R^*, \times)$.
  \item Soit $f \colon \intoo{0}{+\infty} \to \R$ \omterm{def:b1:continuity:continuous}{continue} avec
        $f(xy) = f(x) + f(y)$. Montrer que $f(x) = c\ln x$
        \emph{(transporter par $\exp$)}.
  \item Soit $f \colon \intoo{0}{+\infty} \to \intoo{0}{+\infty}$
        \omterm{def:b1:continuity:continuous}{continue} avec $f(xy) = f(x)f(y)$. Montrer que $f(x) = x^c$.
  \item Trouver toutes les fonctions \omterm{def:b1:continuity:continuous}{continues} $f \colon \R \to \R$
        telles que
        \[
        f(x + y) = f(x) + f(y) + f(x)f(y) .
        \]
        \emph{(Étudier $h = 1 + f$~; traiter à part le cas
        dégénéré.)}
  \item (Parallélogramme) Trouver toutes les fonctions \omterm{def:b1:continuity:continuous}{continues}
        $f \colon \R \to \R$ telles que $f(x + y) + f(x - y) =
        2f(x) + 2f(y)$~: montrer que $f$ est paire, que $f(0) = 0$,
        que $f(nx) = n^2 f(x)$ par récurrence, puis que $f(x) =
        f(1)\,x^2$. (Cette équation est la signature des formes
        quadratiques --- la loi du parallélogramme qui détecte,
        dans le volume de Licence 2, quelles normes proviennent
        d'un produit scalaire.)
\end{enumerate}

\medskip
\noindent\textbf{Partie V --- Jensen et la convexité au milieu.}
\begin{enumerate}[resume]
  \item (Équation de Jensen) Soit $f \colon \R \to \R$ \omterm{def:b1:continuity:continuous}{continue}
        avec $f\bigl(\frac{x+y}{2}\bigr) = \frac{f(x)
        + f(y)}{2}$. Montrer que $g = f - f(0)$ vérifie
        $g\bigl(\frac x2\bigr) = \frac{g(x)}{2}$, en déduire que
        $g$ est additive, et conclure que $f(x) = cx + d$.
  \item Supposons maintenant seulement l'\emph{inégalité}~: $f$
        \omterm{def:b1:continuity:continuous}{continue} avec
        \[
        f\Bigl(\frac{x + y}{2}\Bigr) \leq \frac{f(x) +
        f(y)}{2} \qquad (x, y \in \R).
        \]
        Démontrer par récurrence sur $n$ que, pour tout poids
        dyadique $\lambda = \frac{k}{2^n} \in \intcc{0}{1}$~:
        \[
        f\bigl(\lambda x + (1 - \lambda)y\bigr) \leq \lambda
        f(x) + (1 - \lambda) f(y) .
        \]
  \item Étendre par \omterm{def:b1:continuity:continuous}{continuité} et \omterm{def:b1:topology:dense}{densité} des dyadiques
        (l'\cref{exo:b1:reals:8}) à tout $\lambda \in
        \intcc{0}{1}$~: une fonction \omterm{def:b1:continuity:continuous}{continue} convexe au milieu
        vérifie l'inégalité de convexité complète (la notion
        étudiée systématiquement au \cref{ch:b1:derivative}).
  \item Montrer qu'on ne peut renoncer à la \omterm{def:b1:continuity:continuous}{continuité}~: une
        fonction additive non linéaire $f$ vérifie
        l'\emph{égalité au milieu} de la question 19 sans vérifier
        aucune inégalité de convexité sur un \omterm{prop:b1:reals:intervals}{intervalle} (question
        12). Morale~: la convexité au milieu est une propriété à
        étages dénombrables (les dyadiques), la convexité une
        propriété du \omterm{def:b1:continuity:continuous}{continu}~; la \omterm{def:b1:continuity:continuous}{continuité} est le pont ---
        exactement comme dans les parties I et II.
\end{enumerate}

\medskip
\noindent\textbf{Partie VI --- Dernières variations et synthèse.}
\begin{enumerate}[resume]
  \item Trouver toutes les fonctions \omterm{def:b1:continuity:continuous}{continues} $f \colon \R \to \R$
        telles que $f(x + y) = f(x) + f(y) + xy$ \emph{(soustraire
        la solution particulière $\frac{x^2}{2}$)}.
  \item Démontrer~: si $f \colon \R \to \R$ est \omterm{def:b1:continuity:continuous}{continue} et
        additive seulement sur un \emph{\omterm{def:b1:structures:subgroup}{sous-groupe} \omterm{def:b1:topology:dense}{dense}} $G$
        (c'est-à-dire $f(g + g') = f(g) + f(g')$ pour $g, g' \in
        G$), alors $f$ est additive sur $\R$. Plus généralement,
        deux fonctions \omterm{def:b1:continuity:continuous}{continues} qui coïncident sur une partie
        \omterm{def:b1:topology:dense}{dense} de $\R$ sont égales.
  \item Synthèse, une phrase pour chaque point~: (i) énoncer de
        mémoire l'échelle de régularité de la question 8~; (ii)
        expliquer pourquoi «~graphe \omterm{def:b1:topology:dense}{dense} dans le plan~» est la
        bonne image mentale de l'échec de la régularité~; (iii)
        énumérer les cinq fonctions classiques caractérisées dans
        les parties IV et V et la méthode unique qui les a toutes
        attrapées~; (iv) nommer les deux endroits où la \omterm{def:b1:topology:dense}{densité} de
        $\Q$ (ou des dyadiques) dans $\R$ a porté l'argument, et
        l'endroit où elle ne l'a pas pu (question 13).
\end{enumerate}
\end{problem}
